\section{Discussion}

\subsection{Division of Labor between Rule-Based and Reasoning Agents}

The operational results show a division of labor between two forms of agency rather than between agent and non-agent control. Agent-MD completed 15 system--RH simulation states through 414 campaign-agent operations and 120 run--analyze--decide cycles, while only one production state required review and no automated reasoning-agent invocation occurred during formal production. The rule-based campaign agent remained active throughout production: it selected and dispatched actions from approved rules, task dependencies, and persistent campaign state; generated or updated state-specific inputs; invoked scientific computation and analysis tools; evaluated simulation segments; archived accepted states; validated provenance; and advanced the RH sequence. Its ability to operate without repeated reasoning-agent input arose from persistent workflow state and executable policies rather than repeated open-ended model inference.

The reasoning agent has a different action-selection mechanism and role. It supports flexible campaign construction and interprets unresolved review events, whereas the rule-based campaign agent selects permitted production actions from explicit rules and recorded state. The low number of reasoning-agent invocations therefore does not imply low campaign autonomy or an absence of agentic production. Instead, it shows that persistent state and executable policies allowed the rule-based campaign agent to operate over a long campaign while reserving flexible interpretation for the small number of tasks that required it.

Event-triggered review defines when rule-based campaign-agent execution should pause. A numerical counter, file path, or threshold may be internally valid while being inconsistent with the intended meaning of the current campaign state, and continuing under such a mismatch can propagate an error into downstream simulations. Review is triggered when the recorded workflow state cannot be safely resolved by approved rules, rather than by elapsed time, operation count, or periodic model polling. The rule-based campaign agent packages the unresolved state as a review event and evidence bundle, and the reasoning agent proposes a response or modification. Schema and policy validation, restrictions on permitted actions, and conditional human oversight limit the effect of that proposal before a validated response returns to the production agent.

The principal contribution of Agent-MD is this explicit division of labor between a persistent rule-based campaign agent and a reasoning agent, together with a validated return path for proposed actions. The rule-based campaign agent coordinates repeated production through scientific tools and explicit state, while the reasoning agent provides flexible interpretation during campaign construction and unresolved review events. Reasoning-agent involvement can therefore remain infrequent and concentrated on decisions whose incorrect interpretation could alter the simulation state or provenance of the campaign, without reducing the campaign agent's persistent production role.

\subsection{Scope, Limitations, and Generalization}

The present study evaluates the Agent-MD control architecture within one stateful GCMC--MD desorption campaign comprising five montmorillonite systems and three RH states. The evidence includes one authentic production state requiring review and two frozen incident replays, while formal production required no automated reasoning-agent invocation. The conclusions therefore concern the organization of control, state, and evidence in a long-running campaign. Model-level diagnostic performance and direct comparisons with continuous LLM control in terms of cost, latency, or recovery quality were not evaluated.

Several important elements of the workflow remain defined or approved by the researcher. The scientific objective, prescribed RH path, monitored observables, molecular models, convergence policies, review-event categories, permitted responses, and contents of the evidence bundle were established during campaign construction or subsequent workflow development. The reasoning agent did not independently validate the force field, redefine the physical problem, or replace scientific judgment regarding the adequacy of the simulation protocol. Agent-MD should therefore be understood as providing automation within a researcher-approved set of rules and permitted actions. Its reliability depends on the correctness of the approved workflow rules, the completeness of persistent campaign state and provenance, the validity of the monitored observables, and the provenance of the records supplied to the rule-based campaign agent and any subsequent reviewer.

The emergence of AI agents also changes how scientific workflows should be designed. Most existing simulation environments were developed for human operators and commonly rely on graphical interaction, manually interpreted log files, implicit directory conventions, and contextual knowledge that is not represented explicitly in machine-readable form. Directly assigning a reasoning agent to control such a workflow does not remove these ambiguities and may instead place hidden state and poorly defined decisions inside an open-ended reasoning loop. Conversely, assigning every operation to a reasoning agent is unnecessary when many actions can be selected and dispatched more reliably by a rule-based campaign agent through validated scientific tools. Progress toward agent-enabled science therefore requires a gradual redesign of scientific workflows around explicit state, structured tool interfaces, machine-readable provenance, verifiable actions, and defined review triggers, while preserving human-readable records and opportunities for scientific oversight. The relevant transition is not the wholesale replacement of human-operated workflows by autonomous LLMs, but the development of scientific environments in which researchers, rule-based campaign agents, reasoning agents, and scientific computation and analysis tools interact through explicit and auditable control interfaces.

Generalization of Agent-MD to other computational problems requires more than replacing the molecular system or simulation engine. The target workflow should provide persistent machine-readable state, segmented or restartable execution, routine decisions that can be represented by explicit policies, and identifiable dependencies between predecessor and successor states. It should represent not only task status but also how an accepted predecessor state initializes a successor state through state inheritance. States requiring review should be expressible as structured evidence packages, and the permitted responses should remain sufficiently constrained to be validated before execution. These features are present in many multistage computational problems, including adsorption campaigns, enhanced-sampling calculations, parameter sweeps, sequential electronic-structure workflows, and coupled particle- or continuum-based simulations. Each application would nevertheless require its own definitions of state, completion, provenance, evidence, and review. Further evaluation should therefore examine additional production campaigns and authentic incidents, compare alternative LLM and human-review strategies, and directly assess event-triggered and continuous reasoning architectures in terms of cost, latency, recovery quality, and auditability.
